\usepackage[utf8]{inputenc}
\usepackage{amsmath}
\usepackage{amsfonts}
\usepackage{amssymb}
\usepackage{amsthm}

% --- Halaman kosong sebagai separator antar-bab (gaya ITS) ---
% Auto-inject teks pada halaman kosong verso yang di-generate LaTeX
\makeatletter
\def\cleardoublepage{\clearpage\if@twoside
    \ifodd\c@page\else
        \thispagestyle{empty}
        \vspace*{\fill}
        \begin{center}Halaman ini sengaja dikosongkan\end{center}
        \vspace*{\fill}
        \newpage
        \if@twocolumn\hbox{}\newpage\fi
    \fi\fi}
\makeatother

% Command eksplisit untuk sekat antar-bab
\newcommand{\halamankosong}{%
    \newpage
    \thispagestyle{empty}
    \vspace*{\fill}
    \begin{center}Halaman ini sengaja dikosongkan\end{center}
    \vspace*{\fill}
    \cleardoublepage
}

% --- Pemenggalan kata Indonesia dan toleransi margin ---
\emergencystretch=3em
\hyphenation{
  ak-si-si-me-tris
  an-ti-ko-mu-ta-tor
  an-ti-ko-mu-ta-ti-vi-tas
  an-ti-si-me-tri-sa-si
  ber-ke-leng-kung-an
  ber-ko-res-pon-den-si
  ber-trans-for-ma-si
  di-de-kom-po-si-si
  di-des-krip-si-kan
  di-fe-ren-si-a-bel
  di-for-mu-la-si-kan
  di-i-den-ti-fi-ka-si
  di-ka-rak-te-ri-sa-si
  di-re-nor-ma-li-sa-si
  di-re-pre-sen-ta-si-kan
  di-sub-sti-tu-si-kan
  di-trans-for-ma-si-kan
  elek-tro-di-na-mi-ka
  elek-tro-mag-ne-tik
  elek-tro-mag-ne-tis-me
  ekui-va-len
  ekui-va-len-si
  ho-me-o-mor-fis-me
  in-fi-ni-te-si-mal
  ka-rak-te-ris-tik
  ke-leng-kung-an
  kom-pa-ti-bi-li-tas
  ko-mu-ta-ti-vi-tas
  kon-sis-ten-si-nya
  me-for-mu-la-si-kan
  mem-for-mu-la-si-kan
  mem-per-ta-han-kan
  mem-va-ri-a-si-kan
  meng-eks-ploi-ta-si
  meng-eks-plo-ra-si
  meng-i-den-ti-fi-ka-si
  meng-in-te-gra-si-kan
  meng-kon-fir-ma-si
  meng-kon-struk-si
  me-re-kon-struk-si
  me-re-pre-sen-ta-si-kan
  or-to-go-na-li-tas
  pa-ra-me-te-ri-sa-si
  pe-nye-der-ha-na-an
  se-mi-se-der-ha-na
  sub-sti-tu-si-kan
  ter-lo-ka-li-sa-si
  ter-ma-ni-fes-ta-si
  ter-mo-di-na-mi-ka
}
\usepackage{tikz}
\usetikzlibrary{arrows.meta,decorations.pathreplacing,positioning}

% Theorem environments
\newtheorem{definition}{Definisi}[chapter]
\usepackage{graphicx}
\usepackage{subcaption}
\usepackage[bahasa]{babel}
\usepackage[left=3.0cm, right=2.00cm, top=3.0cm, bottom=2.50cm]{geometry}
\usepackage{tocbibind}
\usepackage{multicol}
\usepackage{indentfirst}
\usepackage{multirow}
\usepackage{array}
\usepackage[table]{xcolor}
\usepackage{float}
\usepackage{mathtools}
\usepackage{cancel} % ditambahkan: diperlukan oleh anotasi \cancel pada Bab 4
\usepackage{wrapfig}
\usepackage[titletoc]{appendix}
\usepackage{listings}
\usepackage{physics}
\usepackage[hidelinks]{hyperref}
%\usepackage{hyperref}
\usepackage{cite}
\usepackage[nosectionbib]{apacite}
\usepackage{background}
\usepackage{tabularx}
\usepackage{setspace}
\usepackage{booktabs}
\usepackage{lscape}
\usepackage{comment}
\usepackage[export]{adjustbox}
\usepackage{enumitem}
\usepackage{textcomp}
\usepackage{mleftright}
\usepackage{grffile}

% subfigure setup using subcaption (modern)
\captionsetup[subfigure]{labelformat = simple, labelsep = none}
\renewcommand\thesubfigure{(\alph{subfigure})\hspace{0.4em}}

\setlength{\heavyrulewidth}{1.5pt}
\setlength{\abovetopsep}{4pt}

\definecolor{codegreen}{rgb}{0,0.6,0}
\definecolor{codegray}{rgb}{0.5,0.5,0.5}
\definecolor{codepurple}{rgb}{0.58,0,0.82}
\definecolor{backcolour}{rgb}{1,1,1}

\usepackage{color}
\usepackage{eso-pic}
\usepackage{pdfpages}

\definecolor{mygreen}{RGB}{28,172,0} % color values Red, Green, Blue
\definecolor{mylilas}{RGB}{170,55,241}

\lstdefinestyle{mystyle}{
	backgroundcolor=\color{backcolour}, 
	commentstyle=\color{codepurple},
	keywordstyle=\color{blue},
	%keywordstyle=\color{black},
	numberstyle=\tiny\color{codegray},
	stringstyle=\color{purple},
	%stringstyle=\color{black},
	basicstyle=\ttfamily,
	breakatwhitespace=false,         
	breaklines=true,                 
	captionpos=b,                    
	keepspaces=true,                 
	numbers=none,                    
	numbersep=5pt,                  
	showspaces=false,                
	showstringspaces=false,
	showtabs=false,                  
	tabsize=2
}
\lstset{style=mystyle}

\usepackage{fancyhdr}
\pagestyle{fancy}
\fancyhf{}
\renewcommand{\headrulewidth}{0pt}
\fancyfoot[LE,RO]{\thepage}
%\rfoot{\thepage}
\usepackage{etoolbox}
\patchcmd{\chapter}{\thispagestyle{plain}}{\thispagestyle{fancy}}{}{}

\addto{\captionsbahasa}{\renewcommand{\bibname}{Daftar Pustaka}}

\renewcommand{\appendixname}{Lampiran}
